\documentclass[11pt]{amsart}
\usepackage{amsmath,amssymb,amsthm,booktabs,array,hyperref,mathtools}
\usepackage[margin=1in]{geometry}
\hypersetup{pdftitle={Nonnested extremal targets in a rational radix--Mobius semigroup},pdfauthor={Facundo Firmenich},pdfsubject={Arithmetic dynamics and computer-assisted proof},pdfkeywords={arithmetic dynamics, radix transformations, Mobius maps, Bellman nonnesting}}

\newtheorem{theorem}{Theorem}[section]
\newtheorem{proposition}[theorem]{Proposition}
\newtheorem{corollary}[theorem]{Corollary}
\newtheorem{lemma}[theorem]{Lemma}
\theoremstyle{definition}
\newtheorem{definition}[theorem]{Definition}
\newtheorem{remark}[theorem]{Remark}

\title[Nonnested extremal targets]{Nonnested extremal targets in a rational radix--M\"obius semigroup}
\author{Facundo Firmenich}
\thanks{Centro de Estudios del Sur (CEDESUR), Universidad Nacional Arturo Jauretche, Buenos Aires, Argentina. Current address: Barcelona, Spain. ORCID: 0009-0002-6578-3811. Correspondence: f.firmenich@cedesur.org.}
\date{27 October 2026}

\begin{document}

\begin{abstract}
We study a noncommutative semigroup of rational interval maps generated by the integer-radix transformation $T_B(x)=\{Bx\}$ and the M\"obius-conjugate branch $S_B=H\circ T_B\circ H^{-1}$, where $H(x)=(1+x)^{-1}$.  For the decimal benchmark $B=10$, seed $17/19$, and a finite certified terminal set, we define extremal target sets
\[
P_m(10,\kappa)=\{x:q(x)>\kappa 10^m\},\qquad q(a/b)=b/(b-a).
\]
Exact rational enumeration gives the first-hit times
\[
(\tau_1,\tau_2,\tau_3,\tau_4)=(45,85,120,162)
\]
at $\kappa=11$.  The corresponding first-hit states are not nested Bellman optima: optimally continuing an earliest depth-$m$ state reaches the next target strictly later than the global depth-$(m+1)$ optimum, with delays $2$, $19$, and $6$ actions in three consecutive transitions.  This phenomenon is robust, not threshold-fine-tuned: all four hitting times remain unchanged for
\[
\kappa\in[10.6197,11.643),
\]
and all three nonnesting inequalities persist simultaneously for
\[
\kappa\in[10.6197,11.01).
\]
The results give exact arithmetic counterexamples to current target height and attained target depth as sufficient Bellman variables.  We describe the certified computation, threshold atlas, and evidence boundary, and formulate the resulting hidden-state problem independently of any unresolved depth-five optimum.
\end{abstract}

\maketitle

\section{Introduction}

Integer-radix transformations of the form $x\mapsto\{Bx\}$ are classical objects in the theory of expansions and interval dynamics \cite{Renyi1957Representations,Parry1960On,Dajani2002Ergodic}.  M\"obius and Farey-type maps, their natural extensions, and their symbolic and geodesic codings form another mature line of research \cite{Nakada1981Metrical,Series1985The,Isola2003On,Panti2026There,Schmidt2025First,Dajani2026Inducing}.  Recent work also emphasizes finite rational orbits in generalized continued-fraction systems and rational semigroup actions \cite{Berthe2025Multidimensional,Javaheri2026Transitive}, while new natural-extension constructions continue to broaden the class of projective arithmetic maps \cite{Chen2026Strange,Fox2026Natural,Nakada2026Farey}.  The present work studies a concrete coupling of these structures in which an integer-radix map and a M\"obius-conjugate branch act on primitive rationals, with exact reduction after every action.

Our focus is not the introduction of a new radix map or a new version of Bellman's principle.  Rather, we isolate an exact arithmetic phenomenon that arises from their interaction: the state that first reaches one extremal target need not be the best precursor for the next, deeper target.  In the benchmark system considered here, three consecutive target transitions exhibit this failure of nesting.  The effect is large: states with present height of order one can outperform already attained portals with heights of order $10^3$ or $10^4$.

The principle of optimality requires a state description that retains all future-relevant information \cite{Bellman1952On,Light2023The}.  State augmentation for non-Markovian objectives is well known in control and planning \cite{Brafman2017Specifying,Jones2021A}.  What is specific here is an exact arithmetic realization: target depth and present height fail as sufficient variables, while the full rational state remains Markovian.  A companion line of work identifies additional Farey and $p$-adic coordinates that organize the hidden future geometry; those coordinates are not required for the main theorems of this paper.

The contributions are:
\begin{enumerate}
\item a base-$B$ rational radix--M\"obius semigroup with explicit branch matrices;
\item exact first-hit times and first-hit sets for four extremal target depths in the decimal benchmark;
\item exact nonnesting certificates for three consecutive target transitions;
\item an exact threshold-robustness atlas proving that the phenomenon persists on a positive interval of $\kappa$;
\item a reproducible computer-assisted proof protocol based entirely on primitive integer arithmetic.
\end{enumerate}

The depth-five hitting time is not needed here.  Its current exact interval is $186\le \tau_5\le212$, but none of the claims below depends on closing that interval.

\section{The rational radix--M\"obius semigroup}

Fix an integer $B\ge2$.  Define
\[
T_B(x)=\{Bx\},\qquad H(x)=\frac1{1+x}.
\]
We regard $T_B$ as a self-map of $[0,1]$ (with $T_B(1)=0$) and define the conjugate generator
\[
S_B=H\circ T_B\circ H^{-1}
\]
as a partial transformation with natural domain
\[
D_S=H([0,1])=[1/2,1].
\]
Thus the generated object is a semigroup of partial rational transformations; a word is admissible at a state only when every partial generator occurring in it is defined.  If $x=a/b$ is reduced, then the active branch digits are
\[
d_T=\left\lfloor\frac{Ba}{b}\right\rfloor,
\qquad
 d_S=\left\lfloor\frac{B(b-a)}{a}\right\rfloor.
\]

\begin{proposition}[Branch formulae]
For a reduced rational $a/b$,
\[
T_B(a/b)=\frac{Ba-d_Tb}{b}
\]
before primitive reduction.  Whenever $S_B$ is admissible,
\[
S_B(a/b)=\frac{a}{Bb-(B-1+d_S)a}
\]
before primitive reduction.
The corresponding integer matrices are
\[
M_T(d)=\begin{pmatrix}B&-d\\0&1\end{pmatrix},
\qquad
M_S(d)=\begin{pmatrix}1&0\\-(B-1+d)&B\end{pmatrix},
\]
and both have determinant $B$.
\end{proposition}

The two generators need not commute.  In the decimal benchmark,
\[
T(S(17/19))=1/2,
\qquad
S(T(17/19))=9/14.
\]
Thus the action is genuinely noncommutative.

\subsection{Benchmark and terminal set}

For the remainder of the paper we set $B=10$ and use the seed
\[
x_0=17/19.
\]
We remove the following certified forward-invariant terminal set:
\begin{align*}
E=\{&0,1,1/2,1/3,2/3,3/4,4/5,5/6,7/8,7/9,7/10,\\
&2/11,4/11,7/11,9/11,11/20,7/12,2/13,5/13,6/13,7/13,8/13,11/13\}.
\end{align*}
Direct child auditing shows that every admissible child of every element of $E$ remains in $E$.  The accompanying terminal-set audit records all $23$ states and all $46$ attempted $T/S$ branches (admissible or not), with zero admissible escapes from $E$.

\section{Extremal target sets}

For $x=a/b<1$ reduced, define the extremal height
\[
q(x)=\frac1{1-x}=\frac{b}{b-a}.
\]
For $m\ge1$ and $\kappa>0$, define
\[
P_m(10,\kappa)=\{x:q(x)>\kappa 10^m\}.
\]
Let $\tau_m(\kappa)$ be the earliest time at which a nonterminal trajectory from $17/19$ reaches $P_m(10,\kappa)$.

The strict inequality in the target definition matters at robustness boundaries.  If the threshold equals the largest previous height, that previous state is not yet a hit.

\section{Exact portal chronology}

\begin{theorem}[Exact benchmark first hits]
At $\kappa=11$,
\[
\boxed{(\tau_1,\tau_2,\tau_3,\tau_4)=(45,85,120,162).}
\]
The first-hit sets are:
\begin{center}
\begin{tabular}{ccl}
\toprule
$m$ & $\tau_m$ & first-hit states\\
\midrule
1&45&$128/129$\\
2&85&$1202/1203$\\
3&120&$11156/11157,\ 11532/11533,\ 11642/11643$\\
4&162&$115482/115483,\ 124852/124853$\\
\bottomrule
\end{tabular}
\end{center}
\end{theorem}

\begin{proof}[Computer-assisted proof]
Every admissible $T$ and $S$ child was generated in exact primitive integer arithmetic.  The set $E$ was removed only after its forward invariance had been independently checked.  States were deduplicated exactly within each layer.  At depth $162$ the exact layer contains $329{,}106{,}757$ states.  The maximum heights of all preceding layers are below the corresponding target thresholds, while the displayed states exceed them at the stated first-hit layers.  The layer maxima and counts are included in the supplementary threshold atlas.
\end{proof}

\section{Nonnested first-hit states}

For a state $x$ reached at time $t$, let $\sigma_m(x)$ denote the least number of additional actions required to reach $P_m$, with $\infty$ if no such continuation exists.

\begin{theorem}[Three consecutive failures of nesting]
At $\kappa=11$:
\begin{align*}
45+\sigma_2(128/129)&=87>85=\tau_2,\\
85+\sigma_3(1202/1203)&=139>120=\tau_3,\\
120+\sigma_4(11156/11157)&=168>162=\tau_4.
\end{align*}
Moreover, the other two first-hit states at depth three, $11532/11533$ and $11642/11643$, have no depth-four hit by time $162$.
\end{theorem}

\begin{proof}
Independent exact target searches give
\[
128/129\xrightarrow[42\text{ actions}]{}1100/1101,
\]
\[
1202/1203\xrightarrow[54\text{ actions}]{}13542/13543,
\]
and
\[
11156/11157\xrightarrow[48\text{ actions}]{}122088/122089.
\]
The searches enumerate $57{,}470$, $150{,}522$, and $113{,}274$ distinct states, respectively.  Independent no-hit certificates through time $162$ for the remaining depth-three first-hit states rule out a global tie.
\end{proof}

\begin{corollary}
Attained target depth is not a Bellman-sufficient state variable.
\end{corollary}

\begin{corollary}
Current height $q(x)$ is not a Bellman value function.
\end{corollary}

Indeed, immediately before each global target hit, there exist much smaller-height states whose future value is larger.  Their heights are
\[
q(524/679)=679/155<5,
\]
\[
q(6897/9503)=9503/2606<4,
\]
and
\[
q(3277/5599)=5599/2322<3.
\]

\section{All depth-three first-hit states fail to extend optimally}

The first-hit set for $P_3$ has three states.  Exact continuation profiles show:
\begin{center}
\begin{tabular}{lcc}
\toprule
source at time $120$ & arrival by time $162$ & best certified status\\
\midrule
$11156/11157$&no&first $P_4$ hit at global time $168$\\
$11532/11533$&no&no $P_4$ hit through global time $162$\\
$11642/11643$&no&no $P_4$ hit through global time $162$\\
\bottomrule
\end{tabular}
\end{center}
Therefore every depth-three first-hit optimum fails to be a depth-four optimum.

\section{Threshold robustness}

For a state with height $M$ at time $d$, it enters $P_m(10,\kappa)$ precisely when
\[
\kappa<\frac{M}{10^m}.
\]
Thus $\tau_m(\kappa)$ is piecewise constant at finite horizon: it can change only when $\kappa=M_d/10^m$ for some exact layer maximum $M_d$.

\begin{lemma}[Exact robustness interval for a fixed first hit]
Suppose the first depth-$m$ hit at parameter $\kappa_0$ occurs at time $d_*$ with hit height $M_{d_*}$ and previous maximum $M^-$.  Then
\[
\tau_m(\kappa)=d_*
\]
for every
\[
\kappa\in\left[\frac{M^-}{10^m},\frac{M_{d_*}}{10^m}\right).
\]
\end{lemma}

The left endpoint is included because the target inequality is strict.

\begin{theorem}[Four-level chronology is not fine-tuned]
All four exact hitting times
\[
(\tau_1,\tau_2,\tau_3,\tau_4)=(45,85,120,162)
\]
remain unchanged for
\[
\boxed{\kappa\in[10.6197,11.643).}
\]
\end{theorem}

\begin{theorem}[Robust nonnesting interval]
All three nonnesting inequalities remain simultaneously valid for
\[
\boxed{\kappa\in[10.6197,11.01).}
\]
\end{theorem}

\begin{proof}
For each source continuation we combine the exact source earliest-arrival profile with the global record-height profile.  The three inequalities persist on
\[
[10.6197,12.6088),
\]
\[
[2.823093281,11.01),
\]
and
\[
[2.641898181,12.2088),
\]
respectively.  Their intersection is $[10.6197,11.01)$.
\end{proof}

The interval is of positive width and contains the benchmark value $\kappa=11$.

\section{Computational protocol and evidence}

The computation uses exact integer arithmetic only.

\begin{enumerate}
\item Every state is stored as a primitive pair $(a,b)$ with $\gcd(a,b)=1$ and $0\le a\le b$.
\item Every admissible child is generated from the exact branch formulae.
\item Target membership is checked before terminal and viability pruning.
\item The finite terminal set is pruned only after independent forward-invariance certification.
\item Layer states are deduplicated exactly by integer pairs.
\item First-hit layers are certified by the complete absence of earlier hits and the explicit presence of hit states at the claimed layer.
\end{enumerate}

The accompanying repository includes:
\begin{itemize}
\item the global layer maxima and counts through depth $162$;
\item exact continuation profiles from the first-hit states;
\item exact threshold intervals;
\item SHA-256 manifests;
\item an independently written checker for the headline claims.
\end{itemize}

\section{Relation to the literature and scope of novelty}

The two ingredients of the construction belong to well-developed theories.  Radix and $\beta$-transformations originate in the work of R\'enyi and Parry and have a large ergodic literature \cite{Renyi1957Representations,Parry1960On,Dajani2002Ergodic}.  Continued-fraction, Farey, and projective M\"obius maps admit natural extensions, geodesic codings, and transfer-operator descriptions \cite{Nakada1981Metrical,Series1985The,Isola2003On,Merriman2023Geodesic,Panti2026There,Schmidt2025First,Dajani2026Inducing,Nakada2026Farey}.  Recent work studies finite rational orbits and transitive rational actions in generalized continued-fraction settings \cite{Berthe2025Multidimensional,Javaheri2026Transitive}; neighboring natural-extension constructions continue to appear for other projective algorithms \cite{Chen2026Strange,Fox2026Natural}.

The Bellman principle and state augmentation for non-Markovian objectives are likewise classical or well-developed \cite{Bellman1952On,Light2023The,Brafman2017Specifying,Jones2021A}.  Survivor sets or trapping in open expanding systems are established objects \cite{Kong2025Open}, and independently checkable certificates are an active topic in formal reachability verification \cite{Chatterjee2025Fixed}.  We therefore do not claim novelty for radix dynamics, projective/Farey methods, state augmentation, trapping, or certification as general concepts.

The claim here is narrower.  After comparison with these nearby literatures, we are not aware of a result for this exact rational radix--M\"obius system that simultaneously provides (i) certified extremal first-hit times, (ii) explicit nonnesting across consecutive extremal targets, and (iii) a positive exact threshold interval on which that nonnesting persists.  This is a provisional priority statement rather than a proof of historical uniqueness; the repository contains a claim-by-claim literature matrix, and specialist review remains part of the pre-submission validation plan.

\section{Discussion}

The examples distinguish three notions that coincide in simpler stopping problems:
\begin{enumerate}
\item being the earliest state to cross the current target;
\item maximizing the current observable height;
\item minimizing the time to a deeper target.
\end{enumerate}
They are different in the present semigroup.

The full rational state is Markovian, but compressed observables can discard decisive arithmetic phase information.  Companion results show that Farey-neighbor and $p$-adic coordinates can select different controls from the same visible rational; those results will be treated separately to avoid conflating the nonnesting theorem with a proposed minimal hidden-state representation.

Several questions remain open.  Does nonnesting persist for infinite families of bases, seeds, and target parameters?  Can one bound the number of target-optimal states?  What arithmetic invariants predict a latent winner?  How do finite modular traps divide the state space into escape and nonescape basins?  These questions are independent of the unresolved exact value of $\tau_5$.

\section{Limitations and evidence boundary}

The exact threshold atlas in this paper is complete through depth $162$ for the benchmark seed, base, and terminal convention.  The robustness theorems are exact on their stated intervals, but no universality in $B$, seed, or $\kappa$ is claimed.  The hidden-state interpretation is supported by exact counterexamples; minimality of any particular augmented coordinate system is not proved here.  The large layer enumeration is internally reproducible and accompanied by independent checking infrastructure, but an external clean-room replication has not yet been completed.  Finally, the priority claim remains provisional pending closed-database citation chaining and external specialist review.

\section{Data and code availability}

The submission package contains the global layer profile, exact continuation profiles, threshold-atlas generator, trap-independent verifiers, and SHA-256 manifests.  A public archival DOI and Software Heritage identifier will be assigned at release; no identifier is claimed before deposition.

\appendix
\section{Exact evidence tables}

The principal nonnesting table is:
\begin{center}
\begin{tabular}{lrrrr}
\toprule
source&source time&global target time&continued arrival&delay\\
\midrule
$128/129$&45&85&87&2\\
$1202/1203$&85&120&139&19\\
$11156/11157$&120&162&168&6\\
\bottomrule
\end{tabular}
\end{center}

The common exact robustness interval is
\[
[10.6197,11.01),
\]
not a rounded confidence interval.  Its endpoints are rational critical values obtained from exact layer maxima.

\begin{thebibliography}{99}

\bibitem{Renyi1957Representations}
A. R\'enyi,
\emph{Representations for real numbers and their ergodic properties},
Acta Math. Acad. Sci. Hungar. (1957), doi:10.1007/BF02020331.

\bibitem{Parry1960On}
W. Parry,
\emph{On the $\beta$-expansions of real numbers},
Acta Math. Acad. Sci. Hungar. (1960), doi:10.1007/BF02020954.

\bibitem{Dajani2002Ergodic}
K. Dajani and C. Kraaikamp,
\emph{Ergodic Theory of Numbers},
Mathematical Association of America, 2002.

\bibitem{Nakada1981Metrical}
H. Nakada,
\emph{Metrical theory for a class of continued fraction transformations and their natural extensions},
Tokyo J. Math. (1981).

\bibitem{Series1985The}
C. Series,
\emph{The modular surface and continued fractions},
J. London Math. Soc. (1985).

\bibitem{Isola2003On}
S. Isola,
\emph{On the spectrum of Farey and Gauss maps},
arXiv:math/0308017 (2003).

\bibitem{Merriman2023Geodesic}
C. Merriman,
\emph{Geodesic flows and slow downs of continued fraction maps},
arXiv:2307.11700 (2023).

\bibitem{Panti2026There}
G. Panti,
\emph{There is only one Farey map},
Ann. Mat. Pura Appl. (2026), doi:10.1007/s10231-026-01693-1; arXiv:2506.02984.

\bibitem{Schmidt2025First}
T. A. Schmidt,
\emph{First return systems for some continued fraction maps},
arXiv:2511.12835 (2025).

\bibitem{Dajani2026Inducing}
K. Dajani, C. Kraaikamp, and S. Sanderson,
\emph{Inducing contractions of the mother of all continued fractions},
J. Number Theory 278 (2026), 816--874, doi:10.1016/j.jnt.2025.05.015.

\bibitem{Berthe2025Multidimensional}
V. Berth\'e, S. Cantrell, J. Lee, and M. Pollicott,
\emph{Multidimensional statistics for finite orbits of generalised continued fractions},
arXiv:2507.11295 (2025).

\bibitem{Javaheri2026Transitive}
M. Javaheri,
\emph{Transitive actions and purely periodic $N$-expansions},
J. Number Theory 283 (2026), doi:10.1016/j.jnt.2025.11.013.

\bibitem{Chen2026Strange}
Y. Chen, K. Dajani, Y. Hu, and C. Kraaikamp,
\emph{A strange continued fraction associated with the Romik map},
arXiv:2605.21098 (2026).

\bibitem{Fox2026Natural}
J. Fox and T. Garrity,
\emph{The natural extension for the triangle map (a multi-dimensional continued fraction) with an internal symmetry from Young conjugation},
arXiv:2606.28014 (2026).

\bibitem{Nakada2026Farey}
H. Nakada, R. Natsui, and J. Thuswaldner,
\emph{Farey graphs and geodesic expansions of complex continued fractions},
arXiv:2603.28468 (2026).

\bibitem{Silverman2007The}
J. H. Silverman,
\emph{The Arithmetic of Dynamical Systems},
Graduate Texts in Mathematics 241, Springer, 2007.

\bibitem{Anashin2009Applied}
V. Anashin and A. Khrennikov,
\emph{Applied Algebraic Dynamics},
De Gruyter, 2009.

\bibitem{Bellman1952On}
R. Bellman,
\emph{On the theory of dynamic programming},
Proc. Natl. Acad. Sci. USA 38 (1952), 716--719.

\bibitem{Light2023The}
B. Light,
\emph{The principle of optimality in dynamic programming: a pedagogical note},
arXiv:2302.08467 (2023).

\bibitem{Brafman2017Specifying}
R. Brafman, G. De Giacomo, and F. Patrizi,
\emph{Specifying non-Markovian rewards in MDPs using LDL on finite traces},
arXiv:1706.08100 (2017).

\bibitem{Jones2021A}
M. Jones and M. Peet,
\emph{A generalization of Bellman's equation with application to path planning, obstacle avoidance and invariant set estimation},
Automatica / arXiv:2006.08175 (2021).

\bibitem{Kong2025Open}
D. Kong, B. Sun, and Z. Wang,
\emph{Open dynamical systems with a moving hole},
arXiv:2505.02336 (2025).

\bibitem{Chatterjee2025Fixed}
K. Chatterjee, T. Quatmann, M. Sch\"affeler, M. Weininger, T. Winkler, and D. Zilken,
\emph{Fixed point certificates for reachability and expected rewards in MDPs},
arXiv:2501.11467 (2025).

\end{thebibliography}

\end{document}
